\documentclass[a4paper,12pt,fleqn]{scrartcl}
\usepackage[T1]{fontenc}
\usepackage{ucs}
\usepackage[utf8x]{inputenc}
\usepackage{ngerman}
\usepackage[ngerman]{babel}
\usepackage{lastpage}
\usepackage[pdftex]{color,graphicx}
\usepackage{listings}
\usepackage{pdflscape}
\usepackage{longtable}
\usepackage[inner=2cm,outer=2cm,top=1cm,bottom=1.5cm,includeheadfoot]{geometry}
\usepackage{fancyhdr}
% \usepackage{url}

%---- Title
\title{Wochenbericht KW 17}
\author{Christoph Gnip}
\date{ 22.4.-28.4.2013 }
%---- End Title

%---- PageLayout
\pagestyle{fancy}

\setlength{\headsep}{10mm}

\usepackage{eso-pic}

%----------------------------------------------------------------------------
% HEADER --------------------------------------------------------------------
%----------------------------------------------------------------------------
\fancyhead[R]{
  \includegraphics[width=100pt,keepaspectratio]{img/amedo2012.png}
}

\fancyhead[C]{ Wochenbericht KW 17 }

\fancyhead[L]{
  \begin{tabular}[b]{l}
  Christoph Gnip\\
  Projekt: RFID-Evolution
  \end{tabular}
}

%Linie oben
\renewcommand{\headrulewidth}{0.5pt}
%----------------------------------------------------------------------------


%----------------------------------------------------------------------------
%----------------------------------------------------------------------------
%----------------------------------------------------------------------------
\fancyfoot[L]{Stand: \today}
\fancyfoot[R]{\thepage{} von \pageref{LastPage}}

% Linie unten
\renewcommand{\footrulewidth}{0.5pt}
%----------------------------------------------------------------------------

% Import Macros  ------------------------------------------------------------
\newcommand\nn{\newline\newline}

%----------------------------------------------------------------------------
% Start the Document --------------------------------------------------------
%----------------------------------------------------------------------------
\begin{document}

\setlength{\headheight}{36pt}

\begin{titlepage}

% make the title
\maketitle

%this has to be called here, after \maketitle
\thispagestyle{empty}

%----include the time table for this week
%Table generated by Excel2LaTeX from sheet 'Tabelle1'
\label{tbl:TTWeek1}
\begin{table}[right]
    \textbf{Anwesenheit}
    \hfill
    \begin{tabular}{|l|l|l|l|l|l|l|}
    \textbf{Mo} & \textbf{Di} & \textbf{Mi} & \textbf{Do} & \textbf{Fr} &
\textbf{Sa} & \textbf{So} \\

    Bo    & Bo    & Bo    & Ge    & Ge    & x     & x \\
    \end{tabular}%
\end{table}%



\end{titlepage}

%- Section 1 -------------------------------------------------------------------
\section[Allgemeines]{Allgemeines}
Arbeiten an der Masterarbeit aufgenommen.
\subsection{Organisation der Arbeit}
\begin{table}[right]
    \renewcommand{\arraystretch}{1.5}
    \begin{tabular}{lp{11cm}}
      Titel der Arbeit:  [DE] & Entwicklung eines Systems zur
Entfernungsabschätzung für Phasen basiertes UHF RFID Tracking durch Verwendung
evolutionärer Berechnungsverfahren \\
      Titel der Arbeit:  [EN] & Development of a Distance Estimation System for
Phase-Based UHF RFID Tracking by Utilizing Methodes of Evolutionary Computation
\\
      Interne Projektbezeichnung: & PRPS-Evo      \\
      Zeitraum: & 22.April-2.September (20 Wochen)
    \end{tabular}
\end{table}
%
Am Donnerstag, den 18.4., fand ein Treffen mit Herrn Prof. Dr. Frank Bärmann statt. Ihm wurde die Masterthesis vorgestellt und er erklärt
sich als Erstbetreuer für die Arbeit einverstanden.
\newline
Die Anmeldung erfolgt am 2.5. Der Termin verspätet sich aufgrund der Öffnungszeiten des Prüfungsamts.


%- Section 2 ----------------------------------------------------------------
\section[Fortschritt]{Projektfortschritt}
In dieser Woche haben die Arbeiten an der Masterarbeit begonnen. Der Rahmen der Arbeit (Muss-, Soll- und Wunschkriterien)
wurde detaillierter im Pflichtenheft abgesteckt. Das fertige Pflichtenheft wird bis Anfang KW 20 ausgearbeitet.

Weiterhin wurden in dieser Woche die in diesem Projekt verwendeten Entwicklungsumgebung(en) aufgesetzt und bedarfsgerecht
installiert. Die Entwicklungsumgebungen, die im Rahmen dieser Arbeit verwendet werden, werden im Pflichtenheft ausführlicher vorgestellt und diskutiert.

Folgende Umgebungen wurden in dieser Woche aufgesetzt:

%- Section 2.1 --------------------------------------------------------------
\subsection{Projektverwaltung}
Für die Verwaltung dieser Arbeit wird das Versionsverwaltungssystem Git verwendet. Es wird mit GitHub \cite{github} webbasierter
Hosting-Dienst verwendet um die Datensicherung zu gewährleisten. Da GitHub im Rahmen einer Abschlussarbeit verwendet wird, erlaubt
es der Betreiber das Repository den öffentlich Zugang zu verweigern. Somit geht das mit der Geheimhaltung konform. Git wird bereits
in anderen Projekten der Amedo GmbH eingesetzt.

%- Section 2.2 --------------------------------------------------------------
\subsection{Entwicklung}
\begin{itemize}
  \item Installation der virtuellen Maschine, die als Produktivumgebung in
      diesem Projekt verwendet wird, als Betriebssystem wird Ubuntu in der Version 12.04 verwendet
  \item Installation aller zur Kompilierung notwendigen Software (u.A. GCC)
  \item Übersetzung und Installation der Shark-Library auf der Linux-Maschine
\end{itemize}

%- Section 2.3 --------------------------------------------------------------
\subsection{Dokumentation}
\begin{itemize}
  \item Für die Erstellung der Dokumentation und der Thesis selbst, wird
      \LaTeX{} verwendet.

\end{itemize}

%- Section 2.4 --------------------------------------------------------------
\subsection{Shark-Library}
\begin{enumerate}
 \item Bereits in dieser Woche ist es gelungen die Beispiel-Programme der Shark-Library
      zu erstellen.
 \item Einarbeitung in den CMA-ES Algorithmus. Der Algorithmus wurde erfolgreich
      in Matlab umgesetzt und erste Versuche (Benchmarks) wurden mit verschiedenen Populationsgrößen durchgeführt.
      Die Ergebnisse des Benchmarks sind hier im Einzelnen nicht vorgestellt, da sie nur der Einarbeitung dienen.
      Ein Auszug aus dem Quellcode ist im \ Anhang~\ref{lst:cmaes-mat-code} gezeigt.
 \item Bereits in dieser Woche ist es gelungen die Beispiel-Programme der Shark-Library zu erstellen.

\end{enumerate}

\subsection{Recherche}
Am Donnerstag wurde im Rahmen des Aufenthalts an der Westf. Hochschule Recherchearbeiten durchgeführt. Dabei lag der
Schwerpunkt auf der Stand der Technik zum den Themen RFID-Tracking und Evolutionäre Verfahren, insb. die Kombintation aus beiden Themenkomplexen.
Ein interessantes Paper \cite{KALMANandSMOOTHING} der gefundenen Literatur beschreibt, wie man mittes Phasen Messung und Kalmann Filtern \cite{Wiki:1} und einem
Glättungsfilter. Die präsentierte Methode war in ihrem Setup ähnlich dem von der Amedo GmbH verwendeten Aufbau.

%- Section 3 -------------------------------------------------------------------
\section[Probleme]{Probleme}
Die Einrichtung der \LaTeX{} Umgebung dauerte länger als erwartet. Es war ein
Tag Arbeit vorgesehen, aber es wurden zwei benötigt.
Unter Windows gab es Stabilitätsprobleme mit dem verwendeten Editor Kile und dem KBibTeX-Tool.
Die beobachteten Instabilitäten traten nicht in der Ubuntu-Umgebung auf. Ein entsprechender Bug-Report wurde eingereicht.

%- Appendix --------------------------------------------------------------------
\begin{appendix}

\newpage
\huge{Anhänge}
\normalsize
\thispagestyle{plain}
\section{Verwendete Abkürzungen}
\dots{}

\newpage
\thispagestyle{plain}
\section{Matlab CMA-ES Code}
\label{lst:cmaes-mat-code}[blablabla]

\newpage
\thispagestyle{plain}
\section{Tabellen}
  \label{tbl:population-compare}

%   \begin{table}[position specifier]
    \begin{table}[h]
      \centering

    \begin{tabular}{lll}
	\hline
	\textbf{Populationsgröße} & \textbf{Ausführungszeit} & \textbf{Funktion}\\\hline
	x & x & x \\
	y & y & y \\

      \end{tabular}
    \caption{Performancevergleich für unterschiedliche Populationsgrößen}
    \end{table}

\end{appendix}


\newpage
%- Bibliography ----------------------------------------------------------------
\bibliographystyle{alpha}
\bibliography{../bib/mathesis_collection1}

% \bibliography{../bib/mathesis_collection1}{}
% \bibliographystyle{IEEEtran}

\end{document}